\section{模型检验}

\subsection{时间步收敛性}
\label{sec:grid}

固定各问已得到的工艺参数，只改变时间步长重新计算关键结果，见表 \ref{tab:grid}。问题一指定位置温度在 $\Delta t\le0.05\,$s 后已稳定；问题二在 $0.025$ 和 $0.0125\,$s 网格上均报告 $79.60\,$cm/min；问题三面积的相对差为 $0.083\%$。因此以 $0.025\,$s 作为全文最终报告网格，可以兼顾边界判定精度与优化计算量。

\begin{table}[htbp]
  \centering
  \caption{关键结果的时间步收敛性}
  \label{tab:grid}
  \small
  \begin{tabular}{crrrcc}
    \toprule
    $\Delta t$/s & 问题一 3 区中点/\dC & 问题二 $v_{\max}$ & 问题三 $A$ & 问题三可行 & 问题四 $J_{\mathrm{sym}}$ \\
    \midrule
    $0.20$   & $131.1306$ & $79.38$ & $408.4327$ & 否 & $0.0228031$ \\
    $0.10$   & $131.1408$ & $79.53$ & $408.6284$ & 否 & $0.0236606$ \\
    $0.05$   & $131.1459$ & $79.58$ & $409.9962$ & 否 & $0.0241783$ \\
    $\bm{0.025}$ & $\bm{131.1484}$ & $\bm{79.60}$ & $\bm{410.6806}$ & 是 & $\bm{0.0248698}$ \\
    $0.0125$ & $131.1497$ & $79.60$ & $411.0229$ & 是 & $0.0247551$ \\
    \bottomrule
  \end{tabular}
\end{table}

表中还说明了三级网格的必要性。同一组问题三参数在 $0.10$ 和 $0.05\,$s 网格上会被判为不可行，而在细网格上满足约束。这不是物理状态改变，而是最优点的峰值裕量极小，离散峰值误差足以改变约束符号。粗网格适合搜索候选区域，最终可行性必须在细网格上重新判断。

\subsection{参数敏感性与工程余量}
\label{sec:sens}

对五个等效换热参数分别作 $\pm5\%$ 单因素扰动，共 10 种情形。图 \ref{fig:sens} 给出问题三面积的变化。回流段参数 $\eta_{\mathrm{ref}}$ 最敏感，$\pm5\%$ 扰动使面积约变化 $\pm47\,\dC\!\cdot\!\mathrm s$，峰值约变化 $\pm1.5\,\dC$；预热段次之。冷却段参数对 $A_L$ 几乎没有影响，因为第三问面积只积分到峰值，冷却发生在积分区间之后。这一结果与模型的因果顺序一致。

\begin{figure}[htbp]
  \centering
  \includegraphics[width=0.78\textwidth]{paper/parameter_sensitivity.pdf}
  \caption{等效换热参数单因素扰动对问题三面积的影响}
  \label{fig:sens}
\end{figure}

参数减小往往使峰值跌破 $240\,\dC$。理论最小面积解在 10 种扰动下保持可行 6 次，工程余量方案保持可行 8 次，说明增加余量确实提高了稳健性，但尚不能替代多工况实验标定和更完整的不确定性分析。

\subsection{优化算法对照}
\label{sec:algo}

在相同种群规模、仿真预算、初始采样和 COBYLA 精修条件下，对差分进化、实数编码遗传算法和粒子群算法各进行一轮交叉检查，并统一在 $0.025\,$s 网格上复核。结果见表 \ref{tab:algo}。

\begin{table}[htbp]
  \centering
  \caption{三种全局搜索算法的同预算对照}
  \label{tab:algo}
  \small
  \begin{tabular}{lrrc}
    \toprule
    算法 & 最优可行面积/(\dC$\cdot$s) & 运行时间/s & 细网格可行 \\
    \midrule
    约束差分进化 $+$ COBYLA & $411.027$ & $183.4$ & 是 \\
    粒子群 $+$ COBYLA & $412.894$ & $177.8$ & 是 \\
    遗传算法 $+$ COBYLA & $424.780$ & $190.4$ & 是 \\
    \bottomrule
  \end{tabular}
\end{table}

三种方法均找到可行解，差分进化给出的面积最小；在其结果上继续高精度精修后得到正文的 $410.6806\,\dC\!\cdot\!\mathrm s$。由于这里只进行了单轮同预算比较，表 \ref{tab:algo} 用于说明主算法选择是经过交叉检查的，不能据此推断差分进化在所有类似问题中都优于另外两种算法。

\subsection{独立模型互校}

焊接区域很薄，中心与表面温差应当较小。为检查 PDE 离散实现，另建立集中参数模型

\begin{equation}
  \frac{\mathrm dT}{\mathrm dt}=k_j[C(x(t))-T]
  \label{eq:lumped}
\end{equation}

并用四阶 Runge--Kutta 独立求解。两种模型在标定工况下的最大温差为 $0.0614\,\dC$，问题一四个指定位置的差异均不超过 $0.02\,\dC$。两条不同的数值路径得到一致结果，既验证了程序实现，也说明对本题的薄焊接区域，厚度方向温差确实很小。
